% --- Mode d'emploi ----------------------------------------------------------
\chapter*{Comment ce livre est fait}
\markboth{Comment ce livre est fait}{Comment ce livre est fait}
\addcontentsline{toc}{chapter}{Comment ce livre est fait}

Tous les chapitres suivent la même route, dans le même ordre. Au bout de deux
ou trois chapitres, vous saurez où regarder sans chercher.

\vspace{0.8em}

\begin{revision}
On rouvre ce que vous savez déjà et dont le chapitre va se servir : une
identité de seconde, une propriété du collège. Jamais un mot de la notion
nouvelle.
\end{revision}

\begin{automatismes}
Une dizaine de calculs qui doivent sortir sans papier. C'est le
chronomètre, pas la réflexion, qui compte ici.
\end{automatismes}

\noindent Vient ensuite \textbf{le cours}, où alternent plusieurs formes :

\begin{definition}[Ce qu'on pose]
Le mot nouveau, écrit une fois pour toutes. On y revient chaque fois qu'un
doute apparaît.
\end{definition}

\begin{theoreme}[Ce qu'on démontre]
Le résultat qui fait travailler la définition, suivi de sa preuve.
\end{theoreme}

\begin{demonstration}
La preuve est écrite en petit, en retrait. Elle se lit après le théorème, ou
plus tard, ou jamais si le temps manque — mais elle est là.
\end{demonstration}

\begin{visualisation}[la figure qui montre l'idée]
Une figure, et à côté la phrase qu'un professeur dirait en la traçant au
tableau. Regardez d'abord le dessin, lisez la légende ensuite.
\end{visualisation}

\begin{exemple}[le calcul déroulé en entier]
Un cas traité ligne à ligne, sans étape sautée, avec les justifications dans
la marge.
\end{exemple}

\begin{methode}{reconnaître le bon outil}
Quand plusieurs théorèmes s'appliquent, ce bloc dit lequel choisir et à quoi
on le reconnaît.
\end{methode}

\begin{erreurclassique}
La faute que l'on corrige chaque année sur les copies, écrite noir sur blanc
pour que vous la voyiez venir.
\end{erreurclassique}

\noindent Puis on passe la main :

\begin{entrainement}[juste après la notion]
Des exercices placés immédiatement après la notion qu'ils font travailler.
Ils sont rangés du plus simple au plus exigeant : entrez par le premier,
arrêtez-vous où vous butez, et revenez-y le lendemain.
\end{entrainement}

\begin{qcm}
Quatre propositions, une seule est exacte. La bonne réponse ne suffit pas :
sachez dire pourquoi les trois autres sont fausses.
\end{qcm}

\begin{vraifaux}
Une affirmation à trancher, et surtout à justifier — par une démonstration si
elle est vraie, par un contre-exemple si elle est fausse.
\end{vraifaux}

\begin{lecturegraphique}
Une courbe, un tableau de variation, un nuage de points : tout ce qui est
demandé se lit sur le document, sans calcul.
\end{lecturegraphique}

\begin{situation}
Un énoncé venu d'ailleurs que des mathématiques, qu'il faut d'abord traduire
avant de pouvoir le résoudre.
\end{situation}

\begin{annales}
Des exercices repris de sujets d'examen de fin de première, jamais réécrits.
Chacun porte son pays, sa série, sa session et son numéro.
\end{annales}

\begin{jecherche}[une situation à démêler]
Un problème plus ouvert, posé \emph{après} le cours, où plusieurs chemins
mènent au résultat.
\end{jecherche}

\begin{jemeteste}
Un questionnaire de fin de chapitre, à faire seul, pour savoir si le chapitre
est acquis avant de passer au suivant.
\end{jemeteste}

\begin{plusloin}
Une question de baccalauréat choisie parce qu'elle ne réclame rien d'autre
que ce chapitre. Vous verrez ainsi, dès cette année, à quoi ressemble
l'exigence de l'épreuve.
\end{plusloin}

\noindent Chaque partie se termine par un \textbf{devoir surveillé} de deux
heures, noté sur vingt, qui porte sur tous ses chapitres.

\begin{notepourlaclasse}
Ce bloc n'apparaît que dans l'édition du professeur : durée conseillée,
difficulté prévisible, variante pour classe à grand effectif.
\end{notepourlaclasse}
